\chapter{Θεωρητικό υπόβαθρο}
\label{ch:theoria}

Το κεφάλαιο αυτό συγκεντρώνει τις έννοιες που απαιτούνται για την κατανόηση των επόμενων.
Ξεκινά από τις γενικές αρχές των \en{blockchains} και των έξυπνων συμβολαίων, συνεχίζει με
το \en{Sui} και τη γλώσσα \en{Move}, στα οποία στηρίζεται η υλοποίηση, και κλείνει με τα
\en{NFTs} και το πρόβλημα της δίκαιης ανταλλαγής, το οποίο η εργασία καλείται να λύσει.

\section{Blockchain: βασικές έννοιες}
\label{sec:blockchain}

\subsection{Κατανεμημένο βιβλίο συναλλαγών}

Ένα \en{blockchain} είναι ένα κατανεμημένο βιβλίο συναλλαγών (\en{distributed ledger}): μια δομή
δεδομένων στην οποία καταγράφονται συναλλαγές και η οποία τηρείται ταυτόχρονα από πολλούς
κόμβους ενός δικτύου, χωρίς κάποιος από αυτούς να έχει προνομιακό ρόλο
\cite{nakamoto2008}. Οι συναλλαγές ομαδοποιούνται σε \en{blocks}, και κάθε \en{block}
περιέχει τη σύνοψη (\en{hash}) του προηγούμενου. Η αλυσίδα που σχηματίζεται έχει την
ιδιότητα ότι η τροποποίηση ενός παλαιότερου \en{block} αλλάζει τη σύνοψή του και, κατ'
επέκταση, ακυρώνει όλα τα επόμενα. Σε συνδυασμό με το γεγονός ότι κάθε κόμβος διατηρεί
δικό του αντίγραφο, η ιδιότητα αυτή καθιστά το ιστορικό πρακτικά αμετάβλητο.

Κάθε συναλλαγή υπογράφεται ψηφιακά από τον αποστολέα της. Η ταυτότητα ενός χρήστη στο
δίκτυο είναι ένα ζεύγος κλειδιών: το δημόσιο κλειδί, από το οποίο παράγεται η διεύθυνσή
του, και το ιδιωτικό, με το οποίο υπογράφει. Δεν υπάρχει μητρώο χρηστών ούτε κωδικοί
πρόσβασης. Όποιος κατέχει το ιδιωτικό κλειδί μιας διεύθυνσης μπορεί να διαθέσει ό,τι
ανήκει σε αυτήν, και κανείς άλλος.

\subsection{Συναίνεση}

Επειδή δεν υπάρχει κεντρική αρχή, οι κόμβοι πρέπει να συμφωνήσουν μεταξύ τους για το ποιες
συναλλαγές είναι έγκυρες και με ποια σειρά εκτελούνται. Ο μηχανισμός με τον οποίο
επιτυγχάνεται αυτή η συμφωνία ονομάζεται πρωτόκολλο συναίνεσης (\en{consensus}). Το
\en{Bitcoin} εισήγαγε το \en{Proof of Work}, στο οποίο το δικαίωμα πρότασης ενός \en{block}
κερδίζεται με υπολογιστική εργασία \cite{nakamoto2008}. Τα νεότερα δίκτυα χρησιμοποιούν
κατά κανόνα \en{Proof of Stake}: οι επικυρωτές (\en{validators}) δεσμεύουν κεφάλαιο ως
εγγύηση καλής συμπεριφοράς, και η συμφωνία επιτυγχάνεται με πρωτόκολλα ανεκτικά σε
βυζαντινά σφάλματα (\en{Byzantine Fault Tolerant}), τα οποία εγγυώνται ορθότητα εφόσον
λιγότερο από το ένα τρίτο της συνολικής δέσμευσης ελέγχεται από κακόβουλους συμμετέχοντες.
Το \en{Sui} ανήκει στη δεύτερη κατηγορία.

\subsection{Έξυπνα συμβόλαια}

Ο όρος έξυπνο συμβόλαιο (\en{smart contract}) προτάθηκε από τον \en{Szabo} για να
περιγράψει συμφωνίες των οποίων οι όροι επιβάλλονται από λογισμικό αντί από τρίτο μέρος
\cite{szabo1997}. Στο πλαίσιο των \en{blockchains}, καθιερώθηκε με το \en{Ethereum}
\cite{buterin2014}: ένα έξυπνο συμβόλαιο είναι πρόγραμμα αποθηκευμένο στην αλυσίδα, το
οποίο εκτελείται από όλους τους κόμβους με ντετερμινιστικό τρόπο όταν το καλεί μια
συναλλαγή. Η εκτέλεση είναι \emph{ατομική}: είτε ολοκληρώνεται και όλες οι αλλαγές που
επέφερε καταγράφονται, είτε αποτυγχάνει και δεν καταγράφεται καμία. Η ατομικότητα αυτή
είναι το θεμέλιο πάνω στο οποίο στηρίζεται κάθε πρωτόκολλο ανταλλαγής χωρίς μεσάζοντα,
όπως θα φανεί στην Ενότητα~\ref{sec:fair-exchange}.

Η εκτέλεση έχει κόστος, το οποίο καταβάλλει ο αποστολέας της συναλλαγής στο εγγενές
νόμισμα του δικτύου και ονομάζεται \en{gas}. Το \en{gas} αποζημιώνει τους επικυρωτές και
ταυτόχρονα αποτρέπει την κατάχρηση των κοινόχρηστων υπολογιστικών πόρων.

\section{Το Sui blockchain}
\label{sec:sui}

Το \en{Sui} είναι ένα δημόσιο \en{blockchain} \en{Proof of Stake} που αναπτύχθηκε από τη
\en{Mysten Labs} και τέθηκε σε λειτουργία το 2023 \cite{sui-whitepaper, sui-docs}. Διαφέρει από τα
περισσότερα δίκτυα της γενιάς του \en{Ethereum} σε δύο θεμελιώδεις επιλογές: η κατάσταση
της αλυσίδας οργανώνεται γύρω από \emph{αντικείμενα} αντί για λογαριασμούς, και τα έξυπνα
συμβόλαια γράφονται στη γλώσσα \en{Move}. Και οι δύο επιλογές επηρεάζουν άμεσα τον
σχεδιασμό της παρούσας εργασίας.

\subsection{Αντικειμενοκεντρικό μοντέλο}
\label{sec:object-model}

Στο \en{Ethereum}, η κατάσταση είναι ένας χάρτης από διευθύνσεις σε λογαριασμούς, και κάθε
συμβόλαιο διαθέτει δικό του ιδιωτικό αποθηκευτικό χώρο. Ένα \en{NFT} του προτύπου
\en{ERC-721} δεν υπάρχει ως αυτοτελής οντότητα· είναι μια εγγραφή της μορφής
\code{tokenId → owner} μέσα στον χώρο του συμβολαίου της συλλογής \cite{erc721}. Η
«μεταβίβαση» είναι μια κλήση που αλλάζει την εγγραφή αυτή.

Στο \en{Sui}, η βασική μονάδα αποθήκευσης είναι το \emph{αντικείμενο} (\en{object})
\cite{sui-object-model}. Κάθε αντικείμενο έχει:
\begin{itemize}[nosep]
  \item ένα μοναδικό αναγνωριστικό (\en{object ID}) 32 \en{bytes},
  \item έναν τύπο, που ορίζεται από κάποιο \en{module} της \en{Move},
  \item έναν αριθμό έκδοσης, που αυξάνεται σε κάθε τροποποίηση,
  \item έναν ιδιοκτήτη, και
  \item τα πεδία δεδομένων που προβλέπει ο τύπος του.
\end{itemize}
Ένα \en{NFT} είναι λοιπόν ένα αντικείμενο πρώτης τάξης. Η διεύθυνση του κατόχου του δεν
βρίσκεται μέσα σε κάποιο συμβόλαιο αλλά στα μεταδεδομένα του ίδιου του αντικειμένου, και η
μεταβίβασή του είναι εγγενής λειτουργία του συστήματος.

Ο ιδιοκτήτης ενός αντικειμένου καθορίζει ποιος μπορεί να το χρησιμοποιήσει ως είσοδο σε
συναλλαγή. Διακρίνονται τέσσερις περιπτώσεις, που απεικονίζονται στο
Σχήμα~\ref{fig:ownership}:

\begin{description}[leftmargin=1.2em,itemsep=2pt]
  \item[Αντικείμενο διεύθυνσης (\en{address-owned}).] Ανήκει σε μία διεύθυνση και μόνον
    αυτή μπορεί να το χρησιμοποιήσει. Έτσι κρατούνται τα \en{NFTs} και τα νομίσματα στα
    πορτοφόλια των χρηστών.
  \item[Κοινόχρηστο αντικείμενο (\en{shared}).] Δεν ανήκει σε κανέναν και μπορεί να το
    χρησιμοποιήσει οποιοσδήποτε, υπό τους κανόνες που επιβάλλει το \en{module} που το
    όρισε. Χρησιμοποιείται όταν πολλοί ανεξάρτητοι χρήστες πρέπει να αλληλεπιδράσουν με την
    ίδια κατάσταση.
  \item[Αμετάβλητο αντικείμενο (\en{immutable}).] Δεν μεταβάλλεται ποτέ ξανά και είναι
    αναγνώσιμο από όλους. Τα δημοσιευμένα πακέτα κώδικα ανήκουν σε αυτή την κατηγορία.
  \item[Αντικείμενο αντικειμένου (\en{object-owned}).] Ανήκει σε άλλο αντικείμενο, για
    παράδειγμα ως δυναμικό πεδίο του. Είναι προσβάσιμο μόνο μέσω του γονέα του.
\end{description}

Στο σύστημα της εργασίας εμφανίζονται και οι τέσσερις κατηγορίες. Ένα \en{NFT} στο
πορτοφόλι ενός χρήστη είναι αντικείμενο με ιδιοκτήτη διεύθυνση. Η προσφορά \en{Bid} είναι
κοινόχρηστο αντικείμενο, ώστε να μπορεί ο προσφέρων να την ακυρώσει και ο κάτοχος του
ζητούμενου \en{NFT} να την αποδεχθεί ή να την απορρίψει. Το δημοσιευμένο πακέτο του
συμβολαίου είναι αμετάβλητο αντικείμενο. Και τα \en{NFTs} που δεσμεύονται μέσα σε μια
προσφορά γίνονται, για όσο διαρκεί η δέσμευση, αντικείμενα που ανήκουν στο \en{Bid}: δεν
εμφανίζονται στο πορτοφόλι κανενός και απελευθερώνονται μόνο όταν η προσφορά κλείσει.

\begin{figure}[ht]
\centering
\begin{tikzpicture}[
  font=\small,
  obj/.style={draw, rounded corners=3pt, minimum width=2.6cm, minimum height=1.05cm, align=center, fill=gray!8},
  who/.style={draw, circle, minimum size=0.9cm, fill=white, font=\footnotesize},
  lbl/.style={font=\footnotesize\itshape, align=center},
  >=Stealth]
  \node[who] (alice) at (0,0) {Α};
  \node[obj, right=1.0cm of alice] (owned) {\en{NFT}\\{\footnotesize\en{address-owned}}};
  \draw[->] (alice) -- (owned);
  \node[lbl, below=0.15cm of owned] {μόνο η Α το χρησιμοποιεί};

  \node[obj, right=2.3cm of owned] (shared) {\en{Bid}\\{\footnotesize\en{shared}}};
  \node[who, above left=0.25cm and 0.9cm of shared] (u1) {Α};
  \node[who, below left=0.25cm and 0.9cm of shared] (u2) {Β};
  \draw[->] (u1) -- (shared);
  \draw[->] (u2) -- (shared);
  \node[obj, right=1.3cm of shared, minimum width=2.3cm] (child) {\en{NFT}\\{\footnotesize\en{object-owned}}};
  \draw[->] (shared) -- node[above, font=\scriptsize] {\en{DOF}} (child);
  \node[lbl, above=0.15cm of shared] {οποιοσδήποτε, με τους\\κανόνες του \en{module}};
  \node[lbl, below=0.15cm of child] {μόνο μέσω του γονέα};
\end{tikzpicture}
\caption{Μορφές ιδιοκτησίας αντικειμένων στο \en{Sui}. Η δεξιά διάταξη είναι ακριβώς
εκείνη που χρησιμοποιεί η εργασία: ένα κοινόχρηστο \en{Bid} που κρατά τα προσφερόμενα
\en{NFTs} ως δυναμικά πεδία αντικειμένων (\en{DOF}).}
\label{fig:ownership}
\end{figure}

Ο Πίνακας~\ref{tab:erc721-vs-sui} συνοψίζει τις διαφορές των δύο μοντέλων ως προς τα
\en{NFTs}.

\begin{table}[H]
\centering
\caption{\en{NFT} στο μοντέλο λογαριασμών και στο αντικειμενοκεντρικό μοντέλο}
\label{tab:erc721-vs-sui}
\small
\begin{tabular}{p{3.3cm}p{5.2cm}p{5.2cm}}
\toprule
 & \en{Ethereum (ERC-721)} & \en{Sui} \\
\midrule
Τι είναι το \en{NFT} & Εγγραφή στον χώρο αποθήκευσης του συμβολαίου της συλλογής & Αυτοτελές αντικείμενο με δικό του \en{ID} \\
Πού καταγράφεται ο κάτοχος & Σε \en{mapping} μέσα στο συμβόλαιο & Στο πεδίο \en{owner} του αντικειμένου \\
Μεταβίβαση & Κλήση συνάρτησης του συμβολαίου & Εγγενής λειτουργία (\code{transfer}) \\
Χρήση από τρίτο συμβόλαιο & Με προηγούμενη έγκριση (\en{approve}) & Το αντικείμενο περνά ως όρισμα, κατά τιμή ή αναφορά \\
«Συλλογή» & Ένα συμβόλαιο & Ένας τύπος \en{Move} \\
\bottomrule
\end{tabular}
\end{table}

Η τέταρτη γραμμή του πίνακα έχει ιδιαίτερη σημασία για τα \en{marketplaces}. Στο
\en{Ethereum}, για να διαπραγματευθεί ένα \en{marketplace} ένα \en{NFT} εκ μέρους του
χρήστη, ο χρήστης πρέπει να του έχει παραχωρήσει διαρκή έγκριση διαχείρισης. Η έγκριση
αυτή παραμένει ενεργή και μετά τη συναλλαγή και αποτελεί γνωστή επιφάνεια επίθεσης. Στο
\en{Sui} δεν υπάρχει αντίστοιχη έννοια: ένα συμβόλαιο αποκτά πρόσβαση σε ένα αντικείμενο
μόνο όταν ο κάτοχός του το δώσει ρητά ως όρισμα στη συγκεκριμένη συναλλαγή.

\subsection{Εκτέλεση συναλλαγών και συναίνεση}
\label{sec:sui-consensus}

Κάθε συναλλαγή στο \en{Sui} δηλώνει ρητά τα αντικείμενα που χρησιμοποιεί ως εισόδους. Αυτό
επιτρέπει στους επικυρωτές να γνωρίζουν εκ των προτέρων ποιες συναλλαγές είναι ανεξάρτητες
μεταξύ τους και να τις εκτελούν παράλληλα. Επιτρέπει επίσης μια διάκριση που δεν υπάρχει
στα περισσότερα \en{blockchains} \cite{suilutris2023}:

\begin{itemize}[nosep]
  \item Οι συναλλαγές που αγγίζουν μόνο αντικείμενα διεύθυνσης δεν ανταγωνίζονται καμία
        άλλη, αφού μόνο ο κάτοχος μπορεί να τις υποβάλει. Γι' αυτές αρκεί, κατά κανόνα, η
        επικύρωση από την απαιτούμενη πλειοψηφία των επικυρωτών, χωρίς να χρειάζεται ολική
        διάταξη ως προς τις υπόλοιπες συναλλαγές του δικτύου.
  \item Οι συναλλαγές που αγγίζουν έστω και ένα κοινόχρηστο αντικείμενο πρέπει να
        διαταχθούν, διότι δύο χρήστες μπορεί να επιχειρήσουν ταυτόχρονα να το
        τροποποιήσουν. Η διάταξη γίνεται από το πρωτόκολλο συναίνεσης, που στην τρέχουσα
        μορφή του δικτύου είναι το \en{Mysticeti} \cite{mysticeti2023}, εξέλιξη των
        \en{Narwhal} και \en{Bullshark} \cite{narwhal2022}.
\end{itemize}

Η συνέπεια για τον σχεδιαστή είναι ότι τα κοινόχρηστα αντικείμενα έχουν κόστος και πρέπει
να χρησιμοποιούνται μόνο όπου είναι αναγκαία. Στην περίπτωσή μας είναι: μια προσφορά
πρέπει να μπορεί να την αποδεχθεί ένας χρήστης και να την ακυρώσει ένας άλλος, άρα δεν
μπορεί να ανήκει αποκλειστικά σε κανέναν από τους δύο. Η συναίνεση προσφέρει επιπλέον
ακριβώς την εγγύηση που χρειάζεται ένα πρωτόκολλο ανταλλαγής: αν η αποδοχή και η ακύρωση
υποβληθούν ταυτόχρονα, μία από τις δύο θα διαταχθεί πρώτη και θα εκτελεστεί, ενώ η άλλη θα
αποτύχει επειδή το αντικείμενο της προσφοράς δεν θα υπάρχει πλέον.

\subsection{Κόμβοι του δικτύου}
\label{sec:nodes}

Οι «κόμβοι» που αναφέρθηκαν μέχρι εδώ δεν είναι όλοι ίδιοι. Στο \en{Sui} διακρίνονται
τρεις ρόλοι με διαφορετικές αρμοδιότητες \cite{sui-full-node}.

\begin{description}[leftmargin=1.2em,itemsep=3pt]
  \item[Επικυρωτές (\en{validators}).] Είναι οι κόμβοι που παράγουν την αλυσίδα. Σε κάθε
    εποχή (\en{epoch}), διάρκειας περίπου 24 ωρών, μια επιτροπή επικυρωτών, καθορισμένη
    από το δεσμευμένο κεφάλαιο, παραλαμβάνει τις συναλλαγές, ελέγχει την εγκυρότητά τους,
    τις εκτελεί, τις διατάσσει μέσω του πρωτοκόλλου συναίνεσης όπου χρειάζεται, και
    υπογράφει τα αποτελέσματα. Οι επικυρωτές ομαδοποιούν επίσης τις εκτελεσμένες
    συναλλαγές σε σημεία ελέγχου (\en{checkpoints}), τα οποία αποτελούν την επίσημη,
    πλήρως διατεταγμένη ιστορία του δικτύου. Δεν εξυπηρετούν ερωτήματα εφαρμογών· ο ρόλος
    τους είναι η επικύρωση, όχι η διάθεση δεδομένων.
  \item[Πλήρεις κόμβοι (\en{full nodes}).] Παρακολουθούν τους επικυρωτές, παραλαμβάνουν
    τα σημεία ελέγχου, επαληθεύουν τις υπογραφές τους και αναπαράγουν τοπικά κάθε
    συναλλαγή, ώστε να διατηρούν δικό τους πλήρες αντίγραφο της κατάστασης και του
    ιστορικού. Δεν συμμετέχουν στη συναίνεση και δεν χρειάζονται άδεια ή δεσμευμένο
    κεφάλαιο: οποιοσδήποτε μπορεί να λειτουργήσει έναν. Είναι οι κόμβοι με τους οποίους
    επικοινωνούν τα πορτοφόλια, οι εξερευνητές και οι αποκεντρωμένες εφαρμογές, μέσω των
    διεπαφών \en{JSON-RPC} και \en{gRPC}: δέχονται νέες συναλλαγές και τις προωθούν στους
    επικυρωτές, και απαντούν σε ερωτήματα όπως «ποια αντικείμενα ανήκουν σε αυτή τη
    διεύθυνση» ή «ποια συμβάντα αυτού του τύπου εκπέμφθηκαν», για τα οποία διατηρούν
    τα αντίστοιχα ευρετήρια. Η \en{Mysten Labs} προσφέρει δημόσιους πλήρεις κόμβους για
    κάθε δίκτυο, και αυτούς χρησιμοποιεί η εφαρμογή της εργασίας.
  \item[\en{Indexers}.] Προαιρετικές υπηρεσίες που διαβάζουν τα σημεία ελέγχου από έναν
    πλήρη κόμβο και τα αποθηκεύουν σε σχεσιακή βάση δεδομένων, ώστε να απαντούν σε
    ερωτήματα που ένας πλήρης κόμβος δεν υποστηρίζει, όπως «όλα τα αντικείμενα ενός τύπου»
    ή συγκεντρωτικά στατιστικά, συνήθως μέσω \en{GraphQL}. Απαιτούν υποδομή και εισάγουν
    μια πρόσθετη πηγή αλήθειας που πρέπει να παραμένει συγχρονισμένη με την αλυσίδα. Η
    εφαρμογή της εργασίας σχεδιάσθηκε ώστε να μην τους χρειάζεται.
\end{description}

Η διάκριση έχει πρακτική σημασία για ό,τι ακολουθεί. Όταν στη συνέχεια λέμε ότι «ο κόμβος
ευρετηριάζει τα συμβάντα» ή ότι «η εφαρμογή ρωτά τον κόμβο», εννοούμε πάντοτε έναν πλήρη
κόμβο. Το τοπικό δίκτυο ανάπτυξης της Ενότητας~\ref{sec:sui-networks} εκκινεί έναν
επικυρωτή και έναν πλήρη κόμβο μέσα στην ίδια διεργασία, οπότε η ίδια διάκριση ισχύει και
εκεί.

\subsection{Programmable Transaction Blocks}
\label{sec:ptb}

Μια συναλλαγή στο \en{Sui} δεν περιορίζεται σε μία κλήση συνάρτησης. Έχει τη μορφή
\en{Programmable Transaction Block (PTB)}: μιας ακολουθίας εντολών που εκτελούνται με τη
σειρά, μοιράζονται εισόδους και μπορούν να χρησιμοποιούν η μία τα αποτελέσματα της άλλης
\cite{sui-ptb}. Το \en{PTB} εκτελείται ατομικά ως σύνολο. Οι εντολές που χρησιμοποιεί η εργασία είναι:

\begin{description}[leftmargin=1.2em,itemsep=1pt]
  \item[\code{MoveCall}] Κλήση δημόσιας συνάρτησης ενός \en{module}.
  \item[\code{SplitCoins}] Αποκοπή ποσού από ένα νόμισμα σε νέο νόμισμα.
  \item[\code{TransferObjects}] Μεταβίβαση αντικειμένων σε διεύθυνση.
  \item[\code{MakeMoveVec}] Κατασκευή διανύσματος (\code{vector<T>}) από επιμέρους
    αντικείμενα, ώστε να δοθεί ως όρισμα σε συνάρτηση.
  \item[\code{Publish}] Δημοσίευση πακέτου.
\end{description}

Το Σχήμα~\ref{fig:ptb} δείχνει το \en{PTB} με το οποίο η εφαρμογή υποβάλλει μια προσφορά:
κατασκευάζεται ένα διάνυσμα από τα προσφερόμενα \en{NFTs}, αποκόπτεται από το νόμισμα του
\en{gas} το προσφερόμενο ποσό \en{SUI}, και τα δύο αποτελέσματα δίνονται ως ορίσματα στη
συνάρτηση \code{place\_bid}. Ο χρήστης υπογράφει μία φορά.

\begin{figure}[ht]
\centering
\begin{tikzpicture}[
  font=\small,
  cmd/.style={draw, rounded corners=3pt, minimum width=3.9cm, minimum height=1.25cm, align=center, fill=gray!8},
  inp/.style={draw, rounded corners=2pt, minimum height=0.6cm, inner xsep=5pt, font=\footnotesize, fill=white},
  >=Stealth]
  \node[cmd] (vec) at (0,0) {\textbf{0:} \code{MakeMoveVec}\\{\footnotesize\code{vector<NFT>}}};
  \node[cmd, right=0.7cm of vec] (split) {\textbf{1:} \code{SplitCoins}\\{\footnotesize από το νόμισμα \en{gas}}};
  \node[cmd, right=0.7cm of split, minimum width=4.2cm] (call) {\textbf{2:} \code{MoveCall}\\{\footnotesize\code{place\_bid<O, R>}}};
  \node[inp, above=0.8cm of vec] (nfts) {\en{NFT}$_1$, \en{NFT}$_2$};
  \node[inp, above=0.8cm of split] (amt) {ποσό σε \en{MIST}};
  \node[inp, above=0.8cm of call] (args) {\code{requested\_count}, \code{target\_nft\_id}};
  \draw[->] (nfts) -- (vec);
  \draw[->] (amt) -- (split);
  \draw[->] (args) -- (call);
  \draw[->] (vec.south) -- ++(0,-0.95) -| ([xshift=-0.9cm]call.south);
  \draw[->] (split.south) -- ++(0,-0.45) -| ([xshift=0.4cm]call.south);
  \node[font=\scriptsize, anchor=north west] at ([yshift=-0.95cm]vec.south) {αποτέλεσμα 0};
  \node[font=\scriptsize, anchor=south west] at ([yshift=-0.45cm]split.south) {αποτέλεσμα 1};
\end{tikzpicture}
\caption{Το \en{PTB} της υποβολής προσφοράς: τρεις εντολές, μία υπογραφή, ατομική εκτέλεση.}
\label{fig:ptb}
\end{figure}

\subsection{Gas και αποθήκευση}
\label{sec:gas}

Το κόστος μιας συναλλαγής στο \en{Sui} έχει δύο σκέλη \cite{sui-gas}. Το \emph{κόστος
υπολογισμού} εξαρτάται από τις πράξεις που εκτελούνται και από την τιμή αναφοράς του
\en{gas} της τρέχουσας εποχής. Το \emph{κόστος αποθήκευσης} είναι ανάλογο του μεγέθους των
αντικειμένων που δημιουργούνται ή μεγαλώνουν και λειτουργεί ως εγγύηση: όταν ένα
αντικείμενο διαγραφεί, το μεγαλύτερο μέρος του ποσού που είχε καταβληθεί για την αποθήκευσή
του (99\% με τις τρέχουσες παραμέτρους του δικτύου) επιστρέφεται (\en{storage rebate}) σε
όποιον υποβάλλει τη συναλλαγή της διαγραφής, όχι κατ' ανάγκην σε όποιον το είχε
δημιουργήσει.
Όλα τα ποσά εκφράζονται σε \en{MIST}, με 1~\en{SUI} $= 10^{9}$~\en{MIST}. Ο μηχανισμός
της επιστροφής θα μας απασχολήσει στην αξιολόγηση (Κεφάλαιο~\ref{ch:axiologisi}), διότι
μια προσφορά είναι ακριβώς ένα αντικείμενο που δημιουργείται για να διαγραφεί.

\subsection{Συμβάντα}
\label{sec:events}

Ένα \en{module} μπορεί να εκπέμπει \emph{συμβάντα} (\en{events}) κατά την εκτέλεση. Τα
συμβάντα δεν αποθηκεύονται στην κατάσταση της αλυσίδας και δεν είναι αναγνώσιμα από άλλα
συμβόλαια· καταγράφονται στα αποτελέσματα της συναλλαγής και ευρετηριάζονται από τους
πλήρεις κόμβους (Ενότητα~\ref{sec:nodes}), ώστε οι εφαρμογές να μπορούν να τα αναζητούν ανά τύπο. Μαζί με κάθε
συμβάν ο κόμβος καταγράφει τον αποστολέα και το ψηφιακό αποτύπωμα της συναλλαγής που το
εξέπεμψε, οπότε ένα ωφέλιμο φορτίο δεν χρειάζεται να επαναλαμβάνει αυτά τα στοιχεία. Η σημασία τους
είναι μεγάλη, επειδή ένας πλήρης κόμβος δεν προσφέρει ερώτημα της μορφής «όλα τα
αντικείμενα τύπου Τ». Μπορεί να απαντήσει ποια αντικείμενα ανήκουν σε μια διεύθυνση, όχι
όμως ποια κοινόχρηστα αντικείμενα ενός τύπου υπάρχουν. Τα συμβάντα είναι επομένως ο
τυπικός τρόπος με τον οποίο μια εφαρμογή ανακαλύπτει τι υπάρχει στην αλυσίδα χωρίς δικό
της \en{indexer}.

\subsection{Πρότυπα του οικοσυστήματος}
\label{sec:sui-standards}

Δύο πρότυπα του \en{Sui} αφορούν άμεσα την εργασία.

\paragraph{Display.} Τα πεδία ενός αντικειμένου είναι αυθαίρετα· ένα πορτοφόλι δεν
γνωρίζει ποιο από αυτά είναι το όνομα ή η εικόνα. Το πρότυπο \en{Object Display}
\cite{sui-display} επιτρέπει στον εκδότη ενός τύπου να ορίσει ένα σύνολο προτύπων
(\en{templates}), π.χ. \code{name → "\{name\}"}, τα οποία οι πλήρεις κόμβοι εφαρμόζουν στα
πεδία του αντικειμένου για να παραγάγουν τα μεταδεδομένα προβολής του. Ο ορισμός απαιτεί
το αντικείμενο \code{Publisher}, που αποδεικνύει ότι ο καλών είναι ο εκδότης του πακέτου.

\paragraph{Wallet Standard.} Το \en{Sui Wallet Standard} \cite{wallet-standard} ορίζει τη
διεπαφή μέσω της οποίας μια εφαρμογή ιστού ανακαλύπτει τα εγκατεστημένα πορτοφόλια και
τους ζητά να υπογράψουν συναλλαγές. Χάρη σε αυτό, η εφαρμογή της εργασίας λειτουργεί με
οποιοδήποτε συμβατό πορτοφόλι χωρίς ειδικό κώδικα για το καθένα.

\subsection{Τα δίκτυα του Sui}
\label{sec:sui-networks}

Το \en{Sui} δεν είναι ένα δίκτυο αλλά τέσσερα, με το ίδιο λογισμικό και διαφορετικό
σκοπό~\cite{sui-networks, sui-local-network}. Η διάκριση έχει σημασία για την εργασία,
διότι καθορίζει πού αναπτύχθηκε, πού μετρήθηκε και πού δημοσιεύεται τελικά το σύστημα
(Ενότητα~\ref{sec:deployment}). Ο Πίνακας~\ref{tab:networks} τα συνοψίζει.

\begin{table}[H]
\centering
\caption{Τα δίκτυα του \en{Sui}}
\label{tab:networks}
\small
\begin{tabular}{l>{\raggedright\arraybackslash}p{3.4cm}>{\raggedright\arraybackslash}p{2.6cm}>{\raggedright\arraybackslash}p{2.7cm}>{\raggedright\arraybackslash}p{3.0cm}}
\toprule
Δίκτυο & Σκοπός & Διατήρηση δεδομένων & Απόκτηση \en{SUI} & Ρόλος στην εργασία \\
\midrule
\en{Mainnet} & Παραγωγή· πραγματικές συναλλαγές & Μόνιμη & Αγορά· πραγματική αξία & Καμία (εκτός εμβέλειας) \\
\en{Testnet} & Δοκιμή πριν την παραγωγή & Μακροχρόνια· σπάνιες, προαναγγελλόμενες διαγραφές & \en{Faucet} με όριο ρυθμού· χωρίς αξία & Τελική δημοσίευση και χρήση με πραγματικό πορτοφόλι \\
\en{Devnet} & Νέες, ασταθείς δυνατότητες & Διαγράφεται εβδομαδιαίως & \en{Faucet} με όριο ρυθμού· χωρίς αξία & Πρώτη δοκιμαστική δημοσίευση \\
\en{Localnet} & Ανάπτυξη στον υπολογιστή του προγραμματιστή & Όσο ζει η διεργασία· νέα γένεση σε κάθε εκκίνηση & Τοπικό \en{faucet} χωρίς όριο & Όλη η ανάπτυξη, οι έλεγχοι και οι μετρήσεις \\
\bottomrule
\end{tabular}
\end{table}

Το \en{mainnet} είναι το δίκτυο παραγωγής: λειτουργείται από ανεξάρτητους επικυρωτές, τα
δεδομένα του διατηρούνται μόνιμα και το \en{SUI} του έχει χρηματική αξία. Το \en{testnet}
τρέχει το ίδιο λογισμικό ως περιβάλλον δοκιμής πριν την παραγωγή· τα νομίσματά του δεν έχουν
αξία και διανέμονται δωρεάν από ένα \en{faucet} με όριο ρυθμού, ενώ τα δεδομένα του
διατηρούνται μακροχρόνια, με σπάνιες διαγραφές που προαναγγέλλονται. Το \en{devnet}
φιλοξενεί τις νεότερες, ενδεχομένως ασταθείς δυνατότητες του πρωτοκόλλου και διαγράφεται
κάθε εβδομάδα, άρα ό,τι δημοσιεύεται εκεί έχει ημερομηνία λήξης. Το \en{localnet}, τέλος,
δεν είναι δημόσιο δίκτυο αλλά μια πλήρης αλυσίδα που εκκινεί ο προγραμματιστής στον δικό
του υπολογιστή με την εντολή \code{sui start}, με έναν επικυρωτή, πλήρη κόμβο και τοπικό
\en{faucet} χωρίς όριο. Κάθε εκκίνηση με την επιλογή \code{-{}-force-regenesis} δημιουργεί
μια νέα, άδεια αλυσίδα με νέο αναγνωριστικό, οπότε κάθε δημοσιευμένο πακέτο και κάθε
αντικείμενο παύουν να υπάρχουν.

Δύο συνέπειες της διάκρισης επανέρχονται στα επόμενα κεφάλαια. Πρώτον, το κόστος σε
\en{gas} είναι ντετερμινιστικό και κοινό σε όλα τα δίκτυα με τις ίδιες παραμέτρους
πρωτοκόλλου, ενώ οι χρόνοι απόκρισης ενός τοπικού δικτύου με έναν επικυρωτή δεν
αντιπροσωπεύουν δημόσιο δίκτυο (Ενότητα~\ref{sec:gas-eval}). Δεύτερον, τα πορτοφόλια
επέκτασης του οικοσυστήματος δηλώνουν ποια δίκτυα υποστηρίζουν, και κατά κανόνα δεν
περιλαμβάνουν το τοπικό· γι' αυτό η ανάπτυξη στο \en{localnet} χρησιμοποιεί ένα
προσωρινό πορτοφόλι μέσα στον φυλλομετρητή, και η δοκιμή με πραγματικό πορτοφόλι γίνεται
στο \en{testnet} (Κεφάλαιο~\ref{ch:dapp}).

\section{Η γλώσσα Move}
\label{sec:move}

Η \en{Move} σχεδιάστηκε αρχικά για το \en{blockchain Diem} με ρητό στόχο την ασφαλή
αναπαράσταση ψηφιακών περιουσιακών στοιχείων \cite{move-paper}. Το \en{Sui} χρησιμοποιεί
μια παραλλαγή της προσαρμοσμένη στο αντικειμενοκεντρικό μοντέλο \cite{move-book}. Ο κώδικας
οργανώνεται σε \emph{modules}, τα οποία ομαδοποιούνται σε \emph{πακέτα} (\en{packages}) και
δημοσιεύονται στην αλυσίδα ως αμετάβλητα αντικείμενα.

\subsection{Πόροι και ικανότητες}
\label{sec:abilities}

Η κεντρική ιδέα της \en{Move} είναι ότι οι τιμές ενός τύπου δεν μπορούν εξ ορισμού ούτε να
αντιγραφούν ούτε να απορριφθούν. Μια τιμή πρέπει να καταλήξει κάπου: να αποθηκευθεί, να
μεταβιβασθεί, ή να αποδομηθεί ρητά από το \en{module} που όρισε τον τύπο της. Ο
μεταγλωττιστής απορρίπτει κάθε πρόγραμμα που αφήνει μια τέτοια τιμή «στον αέρα». Οι
εξαιρέσεις δηλώνονται με \emph{ικανότητες} (\en{abilities}) πάνω στον τύπο, όπως δείχνει ο
Πίνακας~\ref{tab:abilities}.

\begin{table}[H]
\centering
\caption{Οι ικανότητες της \en{Move}}
\label{tab:abilities}
\small
\begin{tabular}{lp{11cm}}
\toprule
Ικανότητα & Τι επιτρέπει \\
\midrule
\code{copy} & Η τιμή μπορεί να αντιγραφεί. \\
\code{drop} & Η τιμή μπορεί να απορριφθεί σιωπηρά στο τέλος της εμβέλειάς της. \\
\code{key} & Η τιμή είναι αντικείμενο του \en{Sui}: έχει πεδίο \code{id: UID} και μπορεί να υπάρχει αυτοτελώς στην αλυσίδα. \\
\code{store} & Η τιμή μπορεί να αποθηκευθεί μέσα σε άλλο αντικείμενο και να μεταβιβασθεί ελεύθερα εκτός του \en{module} που την όρισε. \\
\bottomrule
\end{tabular}
\end{table}

Ένα \en{NFT} δηλώνεται τυπικά ως \code{has key, store}. Η απουσία του \code{copy} σημαίνει
ότι δεν μπορεί να διπλασιασθεί, και η απουσία του \code{drop} ότι δεν μπορεί να χαθεί.
Αυτές δεν είναι συμβάσεις που ο προγραμματιστής οφείλει να τηρεί αλλά ιδιότητες που
ελέγχονται κατά τη μεταγλώττιση και επαληθεύονται εκ νέου από τον \en{bytecode verifier}
κατά τη δημοσίευση. Για ένα πρωτόκολλο δέσμευσης, η συνέπεια είναι ισχυρή: μια συνάρτηση
που παραλαμβάνει ένα διάνυσμα \en{NFTs} δεν μεταγλωττίζεται αν δεν αποφασίσει ρητά πού θα
καταλήξει το καθένα.

\begin{lstlisting}[language=Move,caption={Ένα αντικείμενο με ικανότητες \code{key} και \code{store}},label={lst:abilities}]
public struct Sword has key, store {
    id: UID,
    strength: u64,
}

public fun forge(strength: u64, ctx: &mut TxContext): Sword {
    Sword { id: object::new(ctx), strength }
}

// Δεν μεταγλωττίζεται: το `sword` δεν έχει `drop`
// και δεν μεταβιβάζεται ούτε αποδομείται.
public fun lose(sword: Sword) { }
\end{lstlisting}

\subsection{Γενικοί τύποι}
\label{sec:generics}

Η \en{Move} υποστηρίζει γενικούς τύπους (\en{generics}) με περιορισμούς ικανοτήτων. Η
δήλωση \code{fun wrap<T: key + store>(item: T)} δέχεται οποιονδήποτε τύπο αντικειμένου που
μπορεί να αποθηκευθεί, ανεξάρτητα από το ποιο πακέτο τον όρισε. Οι παράμετροι τύπου
επιλύονται κατά την κλήση, και η εικονική μηχανή ελέγχει ότι τα ορίσματα έχουν πράγματι
τους δηλωμένους τύπους.

Μια παράμετρος τύπου που δεν εμφανίζεται σε κανένα πεδίο της δομής δηλώνεται ως
\code{phantom}. Χρησιμεύει ως ετικέτα σε επίπεδο τύπων: δύο τιμές
\code{Bid<Sword, Shield>} και \code{Bid<Shield, Sword>} έχουν ίδια πεδία αλλά είναι
διαφορετικοί τύποι, και μια συνάρτηση που αναμένει τον έναν δεν δέχεται τον άλλον. Στην
τεχνική αυτή στηρίζεται η υποστήριξη διαφορετικών συλλογών στην εργασία
(Ενότητα~\ref{sec:data-model}).

\subsection{Δυναμικά πεδία}
\label{sec:dynamic-fields}

Τα πεδία μιας δομής ορίζονται στατικά. Όταν ένα αντικείμενο πρέπει να κρατήσει μεταβλητό
πλήθος τιμών, ή τιμές ετερογενών τύπων, χρησιμοποιούνται τα \emph{δυναμικά πεδία}
(\en{dynamic fields}) \cite{sui-dynamic-fields}: ζεύγη ονόματος-τιμής που προσαρτώνται σε
ένα αντικείμενο κατά την εκτέλεση. Υπάρχουν δύο παραλλαγές:

\begin{itemize}[nosep]
  \item Στο απλό δυναμικό πεδίο (\code{dynamic\_field}) η τιμή \emph{περιτυλίσσεται}
        (\en{wrapped}) στο πεδίο. Αν είναι αντικείμενο, παύει να είναι προσβάσιμο με το
        \en{ID} του από εργαλεία εκτός αλυσίδας.
  \item Στο δυναμικό πεδίο αντικειμένου (\code{dynamic\_object\_field}) η τιμή πρέπει να
        είναι αντικείμενο και \emph{παραμένει} αυτοτελές αντικείμενο με το δικό του
        \en{ID}, του οποίου ιδιοκτήτης γίνεται το γονικό αντικείμενο. Οι \en{explorers}
        και τα πορτοφόλια συνεχίζουν να το βλέπουν.
\end{itemize}

Η εργασία χρησιμοποιεί τη δεύτερη παραλλαγή για τα δεσμευμένα \en{NFTs}, ώστε ένα \en{NFT}
που βρίσκεται μέσα σε προσφορά να παραμένει ορατό και επαληθεύσιμο από οποιονδήποτε.

\subsection{Αρχικοποίηση, one-time witness και έλεγχοι}
\label{sec:otw}

Ένα \en{module} μπορεί να ορίσει συνάρτηση \code{init}, που εκτελείται μία μόνο φορά, κατά
τη δημοσίευση του πακέτου. Αν η πρώτη της παράμετρος είναι μια δομή με το όνομα του
\en{module} σε κεφαλαία και μόνη ικανότητα το \code{drop}, το σύστημα κατασκευάζει και της
παραδίδει μια μοναδική τιμή αυτού του τύπου, το λεγόμενο \en{one-time witness (OTW)}.
Επειδή κανείς άλλος δεν μπορεί να κατασκευάσει τέτοια τιμή, το \en{OTW} αποδεικνύει ότι ο
κώδικας εκτελείται κατά τη δημοσίευση, και με αυτό το \en{module} διεκδικεί το αντικείμενο
\code{Publisher} που χρειάζεται για το \en{Display}.

Οι αποτυχίες σηματοδοτούνται με \code{assert!} και έναν αριθμητικό κωδικό (\en{abort code}).
Στην έκδοση 2024 της γλώσσας ο κωδικός δηλώνεται ως σταθερά με την επισημείωση
\code{\#[error]}, η οποία επιτρέπει να συνοδεύεται από περιγραφικό μήνυμα κειμένου που
ενσωματώνεται στο μεταγλωττισμένο πακέτο. Μια αποτυχία ακυρώνει ολόκληρη τη συναλλαγή.
Τέλος, η \en{Move} διαθέτει ενσωματωμένο πλαίσιο ελέγχων μονάδας: το \en{module}
\code{test\_scenario} προσομοιώνει διαδοχικές συναλλαγές από διαφορετικούς αποστολείς,
και η επισημείωση \code{\#[expected\_failure]} επιτρέπει να ελεγχθεί ότι μια κλήση
αποτυγχάνει με τον αναμενόμενο κωδικό και στο αναμενόμενο \en{module}.

\section{NFTs και αγορές}
\label{sec:nfts}

\subsection{Non-Fungible Tokens}

Ένα \en{token} είναι \emph{ανταλλάξιμο} (\en{fungible}) όταν κάθε μονάδα του είναι
ισοδύναμη με κάθε άλλη, όπως συμβαίνει με τα νομίσματα. Ένα \en{Non-Fungible Token} είναι
αντίθετα μοναδικό: αναπαριστά ένα συγκεκριμένο στοιχείο με δική του ταυτότητα
\cite{wang2021nft}. Στο \en{Ethereum} η έννοια τυποποιήθηκε με το \en{ERC-721}
\cite{erc721} και αργότερα με το \en{ERC-1155} \cite{erc1155}, και γνώρισε ευρεία διάδοση
από το 2021 \cite{nadini2021}. Οι χρήσεις περιλαμβάνουν ψηφιακή τέχνη, συλλεκτικά
αντικείμενα, αντικείμενα παιχνιδιών, εισιτήρια και πιστοποιητικά.

Τα \en{NFTs} οργανώνονται σε \emph{συλλογές}: σύνολα αντικειμένων κοινής προέλευσης και
κοινών χαρακτηριστικών. Εντός μιας συλλογής, τα αντικείμενα διαφέρουν συχνά ως προς τη
\emph{σπανιότητα} (\en{rarity}), που επηρεάζει την αξία που τους αποδίδει η αγορά. Στο
\en{Sui}, όπως φάνηκε στον Πίνακα~\ref{tab:erc721-vs-sui}, μια συλλογή αντιστοιχεί
φυσιολογικά σε έναν τύπο \en{Move}. Δύο \en{NFTs} διαφορετικών συλλογών είναι τιμές
διαφορετικών τύπων, και ένα γενικό πρωτόκολλο ανταλλαγής οφείλει να μπορεί να χειρισθεί
τύπους που δεν γνωρίζει εκ των προτέρων.

\subsection{Μηχανισμοί διαπραγμάτευσης}

Οι αγορές \en{NFT} έχουν αναπτύξει μια μικρή οικογένεια μηχανισμών, που συναντώνται με
μικρές παραλλαγές σε όλες τις μεγάλες πλατφόρμες:
\begin{itemize}[nosep]
  \item \emph{καταχώρηση σταθερής τιμής} (\en{listing}): ο κάτοχος ορίζει τιμή και ο πρώτος
        που την καταβάλλει αποκτά το αντικείμενο·
  \item \emph{δημοπρασία} (\en{auction}): οι ενδιαφερόμενοι πλειοδοτούν εντός προθεσμίας·
  \item \emph{προσφορά} (\en{offer} ή \en{bid}): ο ενδιαφερόμενος προτείνει ένα τίμημα για
        ένα συγκεκριμένο αντικείμενο ή για οποιοδήποτε αντικείμενο μιας συλλογής, και ο
        κάτοχος αποφασίζει αν θα το δεχθεί·
  \item \emph{ανταλλαγή} (\en{swap} ή \en{barter}): τα δύο μέρη ανταλλάσσουν αντικείμενα
        μεταξύ τους, ενδεχομένως με συμπληρωματικό ποσό σε νόμισμα.
\end{itemize}
Στους τρεις πρώτους το τίμημα είναι νόμισμα. Η παρούσα εργασία συνδυάζει τον τρίτο με τον
τέταρτο: μια προσφορά της οποίας το τίμημα είναι \en{NFTs}, νόμισμα, ή και τα δύο.

\subsection{Το πρόβλημα της δίκαιης ανταλλαγής}
\label{sec:fair-exchange}

Όταν δύο μέρη που δεν εμπιστεύονται το ένα το άλλο θέλουν να ανταλλάξουν ψηφιακά αγαθά,
ανακύπτει το πρόβλημα της \emph{δίκαιης ανταλλαγής} (\en{fair exchange}): πρέπει είτε να
λάβουν και οι δύο αυτό που αναμένουν, είτε κανείς. Αν ένας στείλει πρώτος, ο άλλος μπορεί
να μην ανταποδώσει. Είναι γνωστό ότι το πρόβλημα δεν λύνεται ντετερμινιστικά χωρίς έμπιστο
τρίτο μέρος \cite{pagnia1999}, και τα κλασικά πρωτόκολλα στηρίζονται σε έναν τέτοιο τρίτο,
που παρεμβαίνει τουλάχιστον σε περίπτωση διαφωνίας \cite{asokan1997}.

Το \en{blockchain} μπορεί να παίξει τον ρόλο αυτού του τρίτου χωρίς να είναι κανείς
συγκεκριμένος. Η ατομικότητα των συναλλαγών εγγυάται ότι οι δύο μεταβιβάσεις είτε
εκτελούνται μαζί είτε καθόλου, ιδέα που έχει μελετηθεί εκτενώς στο πλαίσιο των
\en{atomic swaps} \cite{herlihy2018}. Απομένει ένα πρακτικό ζήτημα: οι δύο πλευρές σπάνια
είναι ταυτόχρονα παρούσες. Ο προσφέρων υποβάλλει την πρότασή του τώρα, και ο αποδέκτης θα
τη δει αργότερα. Στο μεσοδιάστημα, η προσφορά πρέπει να παραμένει εκτελέσιμη.

Υπάρχουν δύο τρόποι να επιτευχθεί αυτό. Στον πρώτο, ο προσφέρων υπογράφει εκτός αλυσίδας
μια εντολή και παραχωρεί στο συμβόλαιο της αγοράς έγκριση να μετακινήσει τα στοιχεία του
όταν βρεθεί αντισυμβαλλόμενος. Τα στοιχεία παραμένουν στο πορτοφόλι του, με συνέπεια η
προσφορά να μπορεί να καταστεί ανεκτέλεστη αν στο μεταξύ τα μετακινήσει. Στον δεύτερο, τα
προσφερόμενα στοιχεία \emph{δεσμεύονται} (\en{escrow}) στην αλυσίδα κατά την υποβολή, υπό
τον έλεγχο του συμβολαίου. Η προσφορά είναι τότε εγγυημένα εκτελέσιμη, με τίμημα ότι τα
στοιχεία δεν είναι διαθέσιμα στον προσφέροντα μέχρι να αποσύρει την προσφορά. Το μοντέλο
αντικειμένων του \en{Sui}, στο οποίο δεν υπάρχουν διαρκείς εγκρίσεις, ευνοεί τον δεύτερο
τρόπο, και αυτόν ακολουθεί η εργασία.
